\section{Introduction}\label{sec:introduction}
A prepared experiment does not supply its observer with a parameter manifold. It supplies legal actions, reports and executable future tests. Past histories are equivalent when those tests cannot distinguish their continuations. Even after this predictive quotient has been realized, its formal size does not determine what a causal observer can retain: the distribution available for compression is created by the observer's own acquisition policy. A finite label subsequently changes that policy and hence changes the next available distribution. The mathematical problem is to characterize this feedback constraint and to determine when its optimal risk can actually be approximated by a finite programme.

There are two distinct obstructions. First, a terminal codebook need not admit intermediate causal updates. Second, separately feasible updates at different phases need not be updates of one stationary programme. Conditional predictive states attached to a reused label may differ with phase as law parameters, but the executed rule may read neither those states nor an uncharged clock. A third issue appears once existence is settled: a Borel report-to-label kernel may require an infinite description. A finite cardinality theorem is not yet an effective finite-resource theorem.

Theorem~\ref{thm:completion} gives an exact allocation invariant and its effective completion. Its first four parts identify feasible retained laws by finite-support barycentric contractions of actual acquisition laws and identify stationary compatibility by one simultaneous affine inequality. The maximum is taken after summing the phase contributions. This eliminates the unknown common report-to-label kernel, but not the law-valued trajectory of occupancies and conditional predictive states. The exact action-slack and pooling-loss identity yields the optimal same-task risk. These exact allocation arguments are retained and re-proved here. Their one-step convex-order ingredient is classical\cite{strassen,leskela}; the serial and stationary compatibility constraints are the additional content.

The effective part is the principal new step. For positive finite-dimensional raw instruments having the support-resolved presentation of Definition~\ref{def:effective}, it yields an interval
\[
 l-E_h-\beta_Q\le B_n^*+\Phi_{\aut}(n,M)\le u,
 \qquad u-l\le\eta,
\]
computed from at most $M2^M K_\rho^M(Q+1)^{MA+M^2\sum_a B_a}$ programmes. Every candidate uses at most the original $M$ labels, including its actions, timing and stopping state. Its rows are shared across phases exactly as the information pattern requires. The error $E_h$ is a complete-law instrument error; $\beta_Q$ separately accounts for probability-row and terminal-decision resolution. The supplied bounds make the interval effectively convergent. In particular, it gives certified lower bounds for \emph{all} Borel programmes, not only upper bounds for a chosen quantization or an enumeration of promising controllers. The external-clock statement is separate and does not charge phase states.

The proof first replaces reports by a common finite partition while retaining their actual cell instruments. A reconstruction measure independent of phase and state averages each old report row into one new common row. Positivity controls the joint physical-state/label error without differentiating a normalized posterior update. A support certificate makes the fine and coarse programmes have the same possible finite paths. Dyadic rounding removes, but never creates, supported edges; thus it retains exact legality and the deadline. Finite enumeration with certified integration then brackets the optimum. Its candidate bound may be very large. It is a value-computation and feasible-execution theorem, not a polynomial-time algorithm or an optimal description-length result.

Theorems~\ref{thm:effective_gaussian} and~\ref{thm:effective_quantum} verify these hypotheses directly for Gaussian hidden-state observations and positive noncommuting instruments. The former permits reducible or deterministic transitions; the latter permits unitary channels and pure states. Neither requires normalized-filter contraction. Proposition~\ref{prop:effective_singular} supplies effective Cantor and atomic-tail refinements and a rank-changing acquired law with its actual masses. These results do not turn the physical hidden or quantum state into a free controller register. The effective presentation, rather than a presumed Fisher or covariance geometry, is checked at the raw kernel.

For accessibility, Theorem~\ref{thm:finite_reuse} solves a complete finite stationary-reuse problem: three labels give risk $1/2$, four give $1/4$, and five give zero. The nontrivial feasible parallelogram and an implementing common kernel are explicit. This example isolates the simultaneous constraint; it is not presented as the resolution of a long-standing filtering problem.

The quantitative geometric theory remains essential. Theorem~\ref{thm:regular} estimates the exact frontier using actual mass, global covering, task curvature and block stability, and gives matching externally timed and internally timed risk curves. Theorem~\ref{thm:blind} gives a different exact curve for noncontracting singular acquisition, controlled by the narrowest intermediate alphabet. These inherited results are re-proved, rather than promoted again as new contributions. Programme bits, scratch, numerical error and parameter-independent simulation appear with explicit certificates in Sections~\ref{sec:effective} and~\ref{sec:resources}. The full-history baseline always belongs to the same horizon, preparation and scored task; a Bayes preparation is not a parameter-uniform minimax theorem.

Finite-state-controller synthesis and observation quantization are established subjects\cite{poupart,junges,continuous}. Our distinction is a certified approximation of the \emph{fixed-cardinality} causal frontier, preserving one whole stationary programme and its exact support constraints. The comparison in Section~\ref{sec:literature} explains both the overlap and the limits. The compact exact theorem alone has no universal effective modulus; the effective theorem assumes known kernels, finite-dimensional positive instruments, support data and computation certificates. Unknown-kernel learning, arbitrary noncompact predictive quotients, infinite horizons and a matched optimum for every computational resource are not claimed. The manuscript's subject remains the general passage from prepared experiments to predictive quotients, actual acquired geometry, executable causal morphisms and finite-resource decision risk.
